\section{ADM split and dynamical closure}
\label{sec:adm_closure}

The variational equations of Section~\ref{sec:variational_derivation} are covariant.
To make contact with protocol time (ticks) and with dynamical initial-value formulations, we summarize the standard $3+1$ (ADM) decomposition.
The key point for CAP-II is not the ADM algebra itself, but where additional constitutive closure is required to interpret lapse and shift in terms of protocol observables.
\par
\noindent In the ADM equations below, $t$ denotes the continuum coordinate time of the effective theory; comparison to protocol logs uses $t=n\tau_0$ after coarse-graining (Section~\ref{sec:introduction} and the conventions in the front matter).

\subsection{Kinematics: lapse, shift, and extrinsic curvature}
\label{subsec:adm_kinematics}

Write the spacetime metric in ADM form:
\begin{equation}
  \dd s^2
  = -\lapse^2\,\dd t^2
  + h_{ij}\left(\dd x^i+N^i\dd t\right)\left(\dd x^j+N^j\dd t\right),
  \label{eq:adm_metric_cap2}
\end{equation}
where $\lapse$ is the lapse, $N^i$ is the shift, and $h_{ij}$ is the induced spatial metric on constant-$t$ slices.
The extrinsic curvature is
\begin{equation}
  K_{ij}
  := \frac{1}{2\lapse}\left(\dot h_{ij}-D_i N_j-D_j N_i\right),
  \label{eq:extrinsic_curvature_cap2}
\end{equation}
where $D_i$ is the covariant derivative of $h_{ij}$ and a dot denotes $\partial_t$.

\subsection{Constraints vs evolution}
\label{subsec:constraints_evolution}

The Einstein equation splits into constraint equations on each slice and evolution equations for $(h_{ij},K_{ij})$.
Schematically, the Hamiltonian constraint and momentum constraint take the form
\begin{align}
  {}^{(3)}R + K^2 - K_{ij}K^{ij}
  &= 16\pi G\,\rho_{\mathrm{tot}} + 2\Lambda, \label{eq:ham_constraint_cap2}\\
  D_j\left(K^{ij}-h^{ij}K\right)
  &= 8\pi G\,j^i_{\mathrm{tot}}, \label{eq:mom_constraint_cap2}
\end{align}
where ${}^{(3)}R$ is the scalar curvature of $h_{ij}$, and $\rho_{\mathrm{tot}},j^i_{\mathrm{tot}}$ are the total energy density and momentum flux seen by the slice normal.

To display the dynamical content explicitly, the evolution equation for the spatial metric is
\begin{equation}
  \dot h_{ij} = -2\lapse K_{ij} + D_i N_j + D_j N_i,
  \label{eq:adm_hdot}
\end{equation}
while the evolution of the extrinsic curvature takes the schematic form
\begin{align}
  \dot K_{ij}
  &=
  -D_iD_j\lapse
  + \lapse\left({}^{(3)}R_{ij}+K K_{ij}-2K_{ik}K^k{}_j\right)
  + \mathcal{L}_{\vec N}K_{ij}
  \nonumber\\
  &\qquad
  -8\pi G\,\lapse\left(S_{ij}-\tfrac12 h_{ij}(S-\rho_{\mathrm{tot}})\right)
  + \Lambda \lapse\,h_{ij},
  \label{eq:adm_Kdot}
\end{align}
where ${}^{(3)}R_{ij}$ is the Ricci tensor of $h_{ij}$, $\mathcal{L}_{\vec N}$ denotes the Lie derivative along the shift vector, and $S_{ij}:=h_i{}^\mu h_j{}^\nu T_{\mu\nu}$ with $S:=h^{ij}S_{ij}$.
See standard ADM references for the full system and conventions \cite{ArnowittDeserMisner1962,Gourgoulhon2007}.

\paragraph{Closure point.}
The protocol-level definition of $\kappa$ directly provides an auditable lapse proxy $\lapse=\kappa_0/\kappa$ (Section~\ref{sec:protocol_dictionary}).
To obtain a quantitative initial-value system one must also specify a spatial coordinate gauge (shift) and identify the slice densities $(\rho_{\mathrm{tot}},j^i_{\mathrm{tot}},S_{ij})$.
In GR the shift is pure gauge, and in the minimal CAP-II closure we adopt a protocol-rest gauge $N^i=0$ (Assumption~\ref{ass:shift_closure}).
The remaining quantities are then determined by the effective stress tensor obtained from the action (Section~\ref{sec:variational_derivation}), with $\kappa$ entering through the $\kappa\mapsto\chi$ identification.
More elaborate realizations may use protocol flow to define a preferred coordinate choice, but the minimal fit targets in Section~\ref{subsec:schwarzschild_fit} do not require this.


